\paragrafos{Μέτρο γωνίας}{Β.1.3, «Μέτρο γωνίας»}

\begin{anagke}
Ο παππούς Γιάννης, που ταξίδευε παλιά με καράβια, φέρνει στην τάξη μια πυξίδα και τη δείχνει στην παρέα. «Οι ναυτικοί βρίσκουν την πορεία τους με την πυξίδα», λέει. «Όταν το καράβι αλλάζει πορεία, ο καπετάνιος δεν λέει “στρίψε λίγο”. Λέει πόσο ακριβώς». Η κυρία Αλεξάνδρα ρωτάει: «Αν το πλοίο στρίψει προς τα βορειοανατολικά ή προς τα βορειοδυτικά, πώς ξέρουμε σε ποια περίπτωση η στροφή είναι μεγαλύτερη;».

Στο μεταξύ στον κήπο η Αγγελική σχεδιάζει μια γωνία φράχτη με πολύ μακριές πλευρές, και ο Νίκος μια ίδια γωνία με κοντές πλευρές. «Η δική μου είναι μεγαλύτερη, έχει μεγαλύτερες πλευρές!», λέει η Αγγελική. «Όχι», απαντά η Ρήνη. «Εγώ θα τις μετρήσω». Παίρνει ένα μοιρογνωμόνιο και μετράει και τις δύο: βγαίνει το ίδιο νούμερο.

Για να συγκρίνουμε, να προσθέτουμε και να αφαιρούμε γωνίες χρειαζόμαστε, όπως και στα μήκη, μια μονάδα και ένα όργανο μέτρησης. Σε αυτή την παράγραφο γνωρίζουμε τη μοίρα και το μοιρογνωμόνιο και μαθαίνουμε να υπολογίζουμε με μέτρα γωνιών.
\end{anagke}

\begin{lexi}{μοίρα, μοιρογνωμόνιο}{μοίρα, μοιρογνωμόνιο}
Η \lexiword{μοίρα} είναι το «μέρος», το μερίδιο που παίρνει ο καθένας. Μία μοίρα είναι ένα πολύ μικρό μέρος του πλήρους κύκλου. Το \lexiword{μοιρογνωμόνιο} είναι το όργανο που «γνωρίζει», δηλαδή μετράει, τις μοίρες.
\end{lexi}

\begin{metafrasi}
Μονάδα μέτρησης των γωνιών. & μοίρα, $1^\circ$\\
Η γωνία $\hat{\mathrm{Α}}$ ανοίγει $40$ μοίρες. & $\hat{\mathrm{Α}}=40^\circ$\\
Υποδιαίρεση της μοίρας. & $1^\circ=60'$ (πρώτα λεπτά)\\
Υποδιαίρεση του λεπτού. & $1'=60''$ (δεύτερα λεπτά)\\
Δύο γωνίες με το ίδιο μέτρο. & $\hat{\mathrm{Κ}}=\hat{\Lambda}$\\
\end{metafrasi}

\tmima{Η θεωρία}

\tmimasmall{Η μοίρα και το μοιρογνωμόνιο}

\begin{orismos}{μοίρα}
\textbf{Μονάδα μέτρησης των γωνιών} είναι η \textbf{μοίρα}, που γράφεται $1^\circ$.
\end{orismos}

Η μέτρηση των γωνιών γίνεται με το \textbf{μοιρογνωμόνιο}. Ο αριθμός που βρίσκουμε ονομάζεται \textbf{μέτρο} της γωνίας. Τοποθετούμε το κέντρο του μοιρογνωμονίου στην κορυφή της γωνίας και τη μία πλευρά πάνω στο μηδέν. Διαβάζουμε το μέτρο στο σημείο όπου η άλλη πλευρά διασχίζει την κλίμακα, και προσέχουμε να διαβάζουμε τη σωστή από τις δύο κλίμακες.

\begin{center}
\begin{tikzpicture}[scale=1.15]
\draw[ggray] (-2.2,0) arc(180:0:2.2) -- cycle;
\foreach \a in {0,10,...,180}{\draw[ggray] ({1.95*cos(\a)},{1.95*sin(\a)})--({2.2*cos(\a)},{2.2*sin(\a)});}
\foreach \a in {0,45,90,135,180}{\node[ggray,font=\scriptsize] at ({1.65*cos(\a)},{1.65*sin(\a)}) {$\a^\circ$};}
\draw[thick] (-2.6,0)--(2.6,0); \draw[thick,gblue] (0,0)--({2.6*cos(40)},{2.6*sin(40)});
\fill (0,0) circle (2pt) node[below] {$\mathrm{Α}$};
\draw[gblue,thick] (0.7,0) arc(0:40:0.7); \node[gblue] at (0.95,0.3) {\small $40^\circ$};
\end{tikzpicture}
\end{center}

Παραδείγματα: το μέτρο της γωνίας $\hat{\mathrm{Α}}$ είναι $40^\circ$ και το μέτρο της $\hat{\mathrm{Β}}$ είναι $125^\circ$. Μια γωνία που ανοίγει όσο η ορθή έχει μέτρο $90^\circ$. Μια γωνία με πλευρές αντικείμενες ημιευθείες έχει μέτρο $180^\circ$ και μια μη κυρτή γωνία έχει μέτρο μεγαλύτερο από $180^\circ$.

Η μοίρα έχει υποδιαίρεση, το \textbf{λεπτό}: $1^\circ=60'$ (πρώτα λεπτά) και $1'=60''$ (δεύτερα λεπτά). Για παράδειγμα, $25^\circ30'=25{,}5^\circ$, γιατί $30'$ είναι η μισή μοίρα. Και $48{,}25^\circ=48^\circ15'$, γιατί $0{,}25\cdot60=15$.

\begin{prosoxi}[Προσοχή: δεν είναι δεκαδικά]
Το $1^\circ=60'$ και όχι $100'$. Γι' αυτό το $25^\circ30'$ δεν είναι $25{,}30^\circ$ αλλά $25{,}5^\circ$.
\end{prosoxi}

\begin{kanonas}{το μέτρο εξαρτάται μόνο από το «άνοιγμα»}
Το μέτρο μιας γωνίας είναι μοναδικό και εξαρτάται μόνο από το \textbf{άνοιγμα} των πλευρών της, όχι από το μήκος τους. Οι πλευρές είναι ημιευθείες και δεν έχουν μήκος.
\end{kanonas}

\tmimasmall{Ίσες γωνίες και σύγκριση γωνιών}

Δύο γωνίες είναι \textbf{ίσες} αν έχουν το ίδιο μέτρο. Τότε η μία εφαρμόζει ακριβώς πάνω στην άλλη.

Μπορούμε να συγκρίνουμε δύο γωνίες και με διαφανές χαρτί, χωρίς μοιρογνωμόνιο:
\begin{itemize}[leftmargin=1.2em,itemsep=0.3ex,topsep=2pt]
\item Αντιγράφουμε τη μία από τις δύο γωνίες σε διαφανές χαρτί.
\item Τοποθετούμε το διαφανές χαρτί πάνω στην άλλη γωνία ώστε να έχουν κοινή κορυφή και μία κοινή πλευρά.
\end{itemize}
Η γωνία που μένει μέσα στην άλλη είναι η μικρότερη.

\begin{center}
\begin{tikzpicture}[scale=1.1]
\begin{scope}
\draw[thick] (3,0)--(0,0)--(2.6,1.5); \draw[gblue] (0.7,0) arc(0:30:0.7); \fill (0,0) circle (2pt);
\node[below left] at (0,0) {$\mathrm{Κ}$}; \node[right] at (1.6,0.3) {\small $30^\circ$};
\end{scope}
\begin{scope}[xshift=4.5cm]
\draw[thick] (1.5,0)--(0,0)--(1.3,0.75); \draw[gblue] (0.7,0) arc(0:30:0.7); \fill (0,0) circle (2pt);
\node[below left] at (0,0) {$\Lambda$};
\end{scope}
\begin{scope}[xshift=8cm]
\draw[thick] (3,0)--(0,0)--(2.2,2); \fill[gblueL] (0,0)--(2.2,0) arc(0:42:2.2)--cycle;
\fill[gwarm] (0,0)--(1.7,0) arc(0:28:1.7)--cycle; \fill (0,0) circle (2pt);
\node[below left] at (0,0) {$\mathrm{Ο}$};
\end{scope}
\end{tikzpicture}
\end{center}

Στα δύο πρώτα σχήματα οι γωνίες έχουν το ίδιο μέτρο ($30^\circ$) παρά τα διαφορετικά μήκη των πλευρών τους, άρα είναι ίσες. Στο τρίτο σχήμα η μικρότερη γωνία (στο εσωτερικό της άλλης) έχει μικρότερο μέτρο.

\tmimasmall{Πρόσθεση και αφαίρεση γωνιών}

Δύο γωνίες μπορούν να προστεθούν με \textbf{επίθεση}: τις τοποθετούμε δίπλα δίπλα, με κοινή κορυφή και μία κοινή πλευρά, χωρίς να επικαλύπτονται. Η γωνία που σχηματίζεται από τις δύο εξωτερικές πλευρές έχει μέτρο το \textbf{άθροισμα} των μέτρων τους. Αντίστροφα, αν από μια γωνία αφαιρέσουμε ένα μέρος της, το υπόλοιπο έχει μέτρο τη \textbf{διαφορά}.

\begin{center}
\begin{tikzpicture}[scale=1.1]
\draw[thick] (3,0)--(0,0); \draw[thick] (0,0)--(2.4,1.7); \draw[thick] (0,0)--(-0.9,2.6);
\fill (0,0) circle (2pt) node[below left] {$\mathrm{Ο}$};
\node[right] at (3,0) {$x$}; \node[right] at (2.4,1.7) {$y$}; \node[above] at (-0.9,2.6) {$z$};
\draw[gblue] (0.8,0) arc(0:35:0.8); \node[gblue] at (1.2,0.18) {\small $\alpha$};
\draw[gwarmD] (0.55,0.38) arc(35:110:0.55) ; \node[gwarmD] at (0.12,0.85) {\small $\beta$};
\end{tikzpicture}
\end{center}

Αν οι γωνίες $x\hat{\mathrm{Ο}}y$ και $y\hat{\mathrm{Ο}}z$ τοποθετηθούν δίπλα δίπλα με αυτό τον τρόπο, τότε $x\hat{\mathrm{Ο}}z=x\hat{\mathrm{Ο}}y+y\hat{\mathrm{Ο}}z$ και $y\hat{\mathrm{Ο}}z=x\hat{\mathrm{Ο}}z-x\hat{\mathrm{Ο}}y$.

\begin{prosoxi}
\begin{itemize}[leftmargin=1.2em,itemsep=0.3ex,topsep=2pt]
\item Μια γωνία με μακρύτερες πλευρές δεν είναι μεγαλύτερη: το μέτρο εξαρτάται μόνο από το άνοιγμα.
\item Στο μοιρογνωμόνιο διαβάζω την κλίμακα που ξεκινά από το μηδέν στην πλευρά που έβαλα πάνω στο $0$. Αν η γωνία δείχνει μικρότερη από την ορθή, δεν γράφω αριθμό μεγαλύτερο από $90$.
\item Για να προσθέσω δύο γωνίες πρέπει να έχουν κοινή κορυφή και κοινή πλευρά και να μην επικαλύπτονται.
\item Το $25^\circ30'$ δεν ισούται με $25{,}30^\circ$.
\end{itemize}
\end{prosoxi}

\tmima{Λυμένα παραδείγματα}

\paradeigma{Η σωστή κλίμακα}
\tixero Το μοιρογνωμόνιο έχει δύο κλίμακες. Η μία πλευρά μιας γωνίας δείχνει στο $0$ και η άλλη δείχνει $40$ στη μία κλίμακα και $140$ στην άλλη. Η γωνία φαίνεται μικρότερη από την ορθή.\par
\tithelo Το μέτρο της γωνίας.\par
\possindeo Αφού η γωνία είναι μικρότερη από την ορθή, το μέτρο της είναι μικρότερο από $90^\circ$.\par
\apantisi $40^\circ$.\par
\elegxos $40^\circ<90^\circ$ και $40^\circ+140^\circ=180^\circ$, όπως οι δύο κλίμακες.

\paradeigma{Μοίρες και λεπτά}
\tixero Η γωνία $\hat{\mathrm{Α}}=36^\circ30'$ και η γωνία $\hat{\mathrm{Β}}=48{,}25^\circ$.\par
\tithelo Να γράψω την $\hat{\mathrm{Α}}$ με δεκαδικό και την $\hat{\mathrm{Β}}$ με λεπτά.\par
\possindeo Επειδή $1^\circ=60'$, τα $30'$ είναι $30:60=0{,}5^\circ$. Και $0{,}25^\circ=0{,}25\cdot60'=15'$.\par
\apantisi $\hat{\mathrm{Α}}=36{,}5^\circ$ και $\hat{\mathrm{Β}}=48^\circ15'$.\par
\elegxos $0{,}5^\circ$ είναι η μισή μοίρα, δηλαδή $30'$.

\paradeigma{Ίσες ή όχι;}
\tixero Η Αγγελική έχει γωνία με πλευρές $12$ cm και ο Νίκος γωνία με πλευρές $4$ cm. Και οι δύο μετρούν $\hat{\omega}=62^\circ$.\par
\tithelo Ποια γωνία είναι μεγαλύτερη;\par
\possindeo Το μέτρο μιας γωνίας εξαρτάται από το άνοιγμα και όχι από τα μήκη των πλευρών.\par
\apantisi Είναι ίσες.\par
\elegxos Αν αντιγράψω τη μία σε διαφανές χαρτί, εφαρμόζει πάνω στην άλλη.

\paradeigma{Πρόσθεση και αφαίρεση}
\tixero Η γωνία του φράχτη είναι $x\hat{\mathrm{Ο}}z=84^\circ$ και ένα μονοπάτι $\mathrm{Οy}$ την χωρίζει ώστε $x\hat{\mathrm{Ο}}y=50^\circ$.\par
\tithelo Η γωνία $y\hat{\mathrm{Ο}}z$.\par
\possindeo Η $x\hat{\mathrm{Ο}}z$ είναι το άθροισμα των $x\hat{\mathrm{Ο}}y$ και $y\hat{\mathrm{Ο}}z$. Άρα $y\hat{\mathrm{Ο}}z=84^\circ-50^\circ$.\par
\apantisi $34^\circ$.\par
\elegxos $50^\circ+34^\circ=84^\circ$.

\paradeigma{Οι δείκτες του ρολογιού}
\tixero Ο κύκλος γύρω από το κέντρο του ρολογιού έχει $360^\circ$ και το ρολόι έχει $12$ ώρες.\par
\tithelo Η γωνία ανάμεσα σε δύο διαδοχικές ώρες και η γωνία των δεικτών στις $4{:}00$.\par
\possindeo Ανάμεσα σε δύο διαδοχικές ώρες: $360^\circ:12=30^\circ$. Στις $4{:}00$ ο λεπτοδείκτης δείχνει το $12$ και ο ωροδείκτης το $4$, δηλαδή $4$ διαστήματα ωρών.\par
\apantisi $30^\circ$ και $4\cdot30^\circ=120^\circ$.\par
\elegxos Στις $3{:}00$ η γωνία είναι $3\cdot30^\circ=90^\circ$, δηλαδή ορθή.

\tmima{Ασκήσεις}

\zoni{Το κατάλαβα}{μέτρο, σύγκριση, λεπτά, πρόσθεση}

\ex Σχεδίασε με το μοιρογνωμόνιο τις γωνίες: $\hat{\alpha}=30^\circ$, $\hat{\beta}=58^\circ$, $\hat{\gamma}=90^\circ$, $\hat{\delta}=135^\circ$, $\hat{\varepsilon}=180^\circ$.

\ex Σχεδίασε δύο γωνίες με μέτρο $40^\circ$: τη μία με πλευρές μήκους $3$ cm και την άλλη με πλευρές μήκους $8$ cm. Ποια είναι μεγαλύτερη; Εξήγησε.

\ex Μετάτρεψε: α) $12^\circ30'$ σε μοίρες με δεκαδικό \quad β) $40^\circ15'$ σε μοίρες με δεκαδικό \quad γ) $18{,}75^\circ$ σε μοίρες και λεπτά \quad δ) $1^\circ$ σε δεύτερα λεπτά.

\ex Βάλε σε αύξουσα σειρά τις γωνίες: $45^\circ$, $40^\circ30'$, $40{,}25^\circ$, $45^\circ10'$.

\ex Οι γωνίες $x\hat{\mathrm{Ο}}y$ και $y\hat{\mathrm{Ο}}z$ έχουν κοινή κορυφή, κοινή πλευρά $\mathrm{Οy}$ και δεν επικαλύπτονται. α) Αν $x\hat{\mathrm{Ο}}y=35^\circ$ και $y\hat{\mathrm{Ο}}z=48^\circ$, βρες την $x\hat{\mathrm{Ο}}z$. β) Αν $x\hat{\mathrm{Ο}}z=120^\circ$ και $x\hat{\mathrm{Ο}}y=75^\circ$, βρες την $y\hat{\mathrm{Ο}}z$.

\ex Σχεδίασε ένα τρίγωνο, μέτρησε τις τρεις γωνίες του με το μοιρογνωμόνιο και βρες το άθροισμά τους. Επανάλαβε με άλλο τρίγωνο. Τι παρατηρείς;

\ex Σωστό ή λάθος; α) Μονάδα μέτρησης των γωνιών είναι η μοίρα. \quad β) Όσο μεγαλύτερο μήκος έχουν οι πλευρές μιας γωνίας, τόσο μεγαλύτερη είναι. \quad γ) Δύο γωνίες είναι ίσες όταν έχουν το ίδιο μέτρο. \quad δ) $1^\circ=100'$. \quad ε) Το μέτρο της ορθής γωνίας είναι $90^\circ$. \quad στ) Το μοιρογνωμόνιο μετράει γωνίες.

\zoni{Το χρησιμοποιώ}{πυξίδα, ρολόι, φράχτης, πόρτα}

\ex \textbf{Το πλοίο.} Ένα πλοίο στρίβει πρώτα $45^\circ$ και μετά ακόμη $30^\circ$ προς την ίδια μεριά. Πόση είναι η συνολική στροφή; Αν μετά στρίψει $20^\circ$ προς την αντίθετη μεριά, πόση στροφή απομένει από την αρχική πορεία;

\ex \textbf{Το ρολόι.} Πόσες μοίρες είναι η γωνία ανάμεσα στους δείκτες στις $2{:}00$, στις $5{:}00$ και στις $6{:}00$;

\ex \textbf{Το παρτέρι.} Η γωνία του φράχτη στον κήπο είναι $84^\circ$. Η Αγγελική σχεδιάζει ένα μονοπάτι που σχηματίζει $50^\circ$ με τη μία πλευρά και χωρίζει το παρτέρι σε δύο μέρη. Πόση γωνία σχηματίζει το μονοπάτι με την άλλη πλευρά;

\ex \textbf{Η πόρτα.} Η πόρτα της τάξης ανοίγει πρώτα $35^\circ$ και μετά ανοίγει ακόμη $50^\circ$. Στη συνέχεια η Φανή τη σπρώχνει και κλείνει $20^\circ$. Πόσο ανοιχτή είναι τώρα;

\ex \textbf{Η κλίμακα.} Η Φανή μετράει μια γωνία που φαίνεται μικρότερη από την ορθή και το μοιρογνωμόνιο δείχνει $130$ στη μία κλίμακα και $50$ στην άλλη. Ποιο είναι το σωστό μέτρο και γιατί;

\zoni{Το σκέφτομαι}{τι εξαρτάται από τι}

\ex Γιατί το μέτρο μιας γωνίας δεν εξαρτάται από το μήκος των πλευρών της; Σκέψου τι είναι οι πλευρές μιας γωνίας.

\ex Οι γωνίες $x\hat{\mathrm{Ο}}y=70^\circ$, $y\hat{\mathrm{Ο}}z=50^\circ$ και $z\hat{\mathrm{Ο}}t=40^\circ$ είναι διαδοχικές, χωρίς να επικαλύπτονται. Πόση είναι η $x\hat{\mathrm{Ο}}t$;

\ex Πώς μπορείς να συγκρίνεις δύο γωνίες όταν δεν έχεις μοιρογνωμόνιο; Περιέγραψε τα βήματα.

\ex Σχεδίασε μια γωνία $x\hat{\mathrm{Ο}}y=30^\circ$. Με κοινή κορυφή και κοινή πλευρά $\mathrm{Οx}$ σχεδίασε γωνία $x\hat{\mathrm{Ο}}t$ α) διπλάσιου μέτρου \quad β) μισού μέτρου.

\zoni{Για το σπίτι}{γωνίες γύρω μας}

\ex Βρες στο σπίτι πέντε γωνίες (ανοιχτό βιβλίο, ψαλίδι, ανοιχτή πόρτα, γωνία τραπεζιού, ανοιχτός φορητός υπολογιστής) και μέτρησε το άνοιγμά τους με το μοιρογνωμόνιο.

\ex Άνοιξε έναν διαβήτη σε τρία διαφορετικά ανοίγματα και σχεδίασε στο τετράδιό σου τη γωνία που σχηματίζουν τα σκέλη του κάθε φορά. Μέτρησέ τες.
